\section{Composition: Multi-Panel Posters}
\label{sec:composition}

The Composition layer arranges multiple grammar outputs into a single poster or infographic. This is how the multi-section corpus images (10x-coding-playbook, archon-harness-builder, rag-vectorless-explainer) are assembled.

\subsection{Design Philosophy}

\paragraph{Composition is a thin layer.}
The Composition layer does not define new visual primitives. It arranges existing grammar outputs (each a Scene IR) into a larger canvas.

\paragraph{Each panel is an independent grammar.}
A panel might be a timeline, a flow diagram, a comparison table, or a stat callout. Each panel has its own Domain IR and is compiled independently.

\paragraph{The composition specifies arrangement, not content.}
Content is defined in each panel's Domain IR. The Composition IR specifies panel positions, sizes, and inter-panel elements (dividers, titles).

\subsection{Composition IR}

\begin{lstlisting}[language=yaml,caption={Composition IR}]
version: "1.0"
metadata:
  title: "The 10x AI Coding Playbook"
  subtitle: "Everything is infrastructure for one thing: parallel development."
  theme: dark-poster
  canvas:
    width: 1200
    height: auto               # computed from panels
composition:
  layout: stack                # stack | grid | custom
  gap: 40                      # px between panels
  panels:
    - id: reframe
      grammar: timeline
      title: "THE REFRAME"
      ir:                      # inline Domain IR
        tracks:
          - id: solo
            label: "2x = solo + AI"
            # ...
    - id: infrastructure
      grammar: flow
      title: "THE INFRASTRUCTURE"
      ir:
        direction: left-to-right
        nodes:
          # ...
    - id: catch
      grammar: flow
      title: "THE CATCH"
      ir:
        # ...
    - id: payoff
      grammar: timeline
      title: "THE PAYOFF"
      ir:
        # ...
  chrome:                      # editorial elements
    header:
      logo: { src: "icon.png", position: top-left }
      title_style: large-centered
    footer:
      badges: ["GitHub", "Linear", "Jira"]
    dividers: between-panels   # none | between-panels | after-each
\end{lstlisting}

\subsection{Layout Modes}

\subsubsection{Stack Layout}

Panels arranged vertically, each taking full width. Most common for long-form infographics.

\begin{Verbatim}[frame=single,fontsize=\small]
+----------------------------------+
|  Panel 1 (full width)            |
+----------------------------------+
|  Panel 2 (full width)            |
+----------------------------------+
|  Panel 3 (full width)            |
+----------------------------------+
\end{Verbatim}

\subsubsection{Grid Layout}

Panels arranged in a grid (rows × columns). Useful for dashboards or comparison layouts.

\begin{lstlisting}[language=yaml,caption={Grid layout specification}]
composition:
  layout: grid
  grid:
    columns: 2
    rows: auto                 # computed from panel count
  panels:
    - { id: p1, span: 1 }      # 1 cell
    - { id: p2, span: 1 }
    - { id: p3, span: 2 }      # 2 cells (full row)
\end{lstlisting}

\subsubsection{Custom Layout}

Explicit (x, y, width, height) for each panel. Maximum control, most verbose.

\begin{lstlisting}[language=yaml,caption={Custom layout specification}]
composition:
  layout: custom
  panels:
    - id: p1
      x: 0
      y: 0
      width: 600
      height: 300
    - id: p2
      x: 600
      y: 0
      width: 600
      height: 300
\end{lstlisting}

\subsection{Compilation Process}

\begin{enumerate}
  \item \textbf{Parse Composition IR}: Validate structure, resolve grammar references.
  
  \item \textbf{Compile each panel}: For each panel, invoke the appropriate grammar's layout engine to produce a Scene IR.
  
  \item \textbf{Compute panel dimensions}: If not explicit, compute from panel Scene dimensions.
  
  \item \textbf{Arrange panels}: Apply layout mode to compute (x, y) for each panel.
  
  \item \textbf{Merge Scenes}: Combine all panel Scenes into a single composition Scene, translating coordinates.
  
  \item \textbf{Add chrome}: Render header, footer, dividers, badges.
  
  \item \textbf{Output}: The final Scene IR (a single scene containing all panels).
\end{enumerate}

\subsection{Panel References}

Panels can reference external IR files instead of inline definitions:

\begin{lstlisting}[language=yaml,caption={Panel with external IR reference}]
panels:
  - id: timeline-panel
    grammar: timeline
    title: "Project Roadmap"
    ir_file: "./roadmap.timeline.yaml"  # external file
\end{lstlisting}

This enables:
\begin{itemize}
  \item Reusing the same diagram in multiple compositions
  \item Separate version control for each panel's IR
  \item Incremental updates (change one panel, not the whole poster)
\end{itemize}

\subsection{Chrome Elements}

\subsubsection{Header}

\begin{itemize}
  \item \textbf{Logo}: positioned at top-left or top-right
  \item \textbf{Title}: large centered text
  \item \textbf{Subtitle}: smaller text below title
\end{itemize}

\subsubsection{Section Titles}

Each panel can have a title rendered above it:

\begin{lstlisting}[language=yaml,caption={Panel title configuration}]
panels:
  - id: infrastructure
    title: "THE INFRASTRUCTURE"
    title_style: section-header   # theme-defined style
\end{lstlisting}

\subsubsection{Dividers}

Horizontal lines between panels:

\begin{lstlisting}[language=yaml,caption={Divider options}]
chrome:
  dividers: between-panels       # none | between-panels | after-each
  divider_style:
    color: "#333"
    thickness: 1
    margin: 20
\end{lstlisting}

\subsubsection{Footer}

\begin{itemize}
  \item \textbf{Badges}: pill-shaped labels (``GitHub'', ``MIT License'')
  \item \textbf{Author}: avatar + handle
  \item \textbf{Branding}: company logo or URL
\end{itemize}

\subsection{Determinism in Composition}

Composition inherits determinism from its constituent grammars:

\begin{itemize}
  \item Each panel's Scene is deterministic (by grammar contract)
  \item Panel arrangement is deterministic (explicit layout)
  \item Chrome rendering is deterministic (explicit configuration)
\end{itemize}

The composition's \texttt{scene\_hash} is computed from the merged Scene, enabling golden testing of complete posters.

\subsection{Limitations}

\begin{description}
  \item[No cross-panel connectors] Edges cannot span from a node in panel A to a node in panel B. Panels are visually independent.
  
  \item[No responsive layout] Panel sizes are computed once; there is no responsive reflow for different canvas widths.
  
  \item[No interactive linking] Clicking a panel does not navigate or filter. Output is static.
\end{description}

These limitations are intentional: they keep the Composition layer simple and deterministic. Cross-panel connections and responsive layout would significantly complicate the architecture.
